\chapter{FPGA 엔진}\label{ch:fpga-engine}

기준 RTL은 칩에 올릴 수 없다(\ref{sec:sim-vs-synth}절). 실제 FPGA에 올리는 것은 같은 동작을 칩에 맞는
구조로 다시 쓴 \emph{FPGA 엔진}(\code{rtl/fpga})이다. 엔진은 이 프로젝트의 세 번째 봄베이기도 하다.
기어로 도는 기계식, 릴레이식과 나란히 놓여 같은 메뉴에서 같은 정지를 내는, 움직이는 부품이 없는 전자식이다.

엔진의 설계는 이 문서를 쓰는 시점에 타이밍을 맞추는 조정이 진행 중이다. 그래서 이 장은 이미 정해진
구조의 개요만 적고, 코드의 해설은 설계가 확정된 뒤에 쓴다.

\section{지금까지 나온 것}

엔진에 대해 1부에서 이미 설명한 것을 모으면 다음과 같다.

\begin{itemize}
\item \textbf{바뀌는 것만 다시 계산한다}(\ref{sec:tradeoff}절). 스크램블러의 치환표를 위치마다 계산하지 않고,
가운데 로터가 바뀔 때에만 계산해 분산 RAM에 넣어 둔다.
\item \textbf{치환표의 계산은 파이프라인이다}(\ref{sec:pipeline}절). 안쪽 치환은 글자를 하나씩 흘려 넣는
14단 파이프라인으로, 오른쪽 로터 26위치의 표는 26글자를 나란히 처리하는 5단으로 계산한다.
\item \textbf{전파는 케이블 단위로 한다}(\ref{ch:lesson}장). 케이블 선택, 치환, 원핫 OR의 세 단계다.
\item \textbf{정지는 FIFO에 넣는다}(\ref{sec:handshake}절).
\item \textbf{엔진 여러 대를 나란히 놓는다.} 엔진 하나가 로터 순서 하나를 맡고, PC나 보드의 처리 장치가
레지스터 버스를 통해 메뉴와 로터 순서를 써넣고 정지를 읽어 간다(그림~\ref{fig:g-engine-blocks}).
\end{itemize}

\begin{figure}[ht]
\centering
\begin{tikzpicture}[x=1cm,y=1cm]
  \node[blkg,minimum width=22mm,minimum height=14mm] (host) at (-0.2,-0.1) {호스트\\{\scriptsize PC 또는 보드의 ARM}};
  \node[blkn,minimum width=22mm,minimum height=14mm] (br) at (2.8,-0.1) {브리지\\{\scriptsize UART / AXI4-Lite}};
  \draw[rounded corners=2pt,draw=cgray,fill=clight] (5.9,-3.1) rectangle (14.5,2.0);
  \node[anchor=north west,font=\small\bfseries] at (6.0,1.95) {\code{bombe\_array}: 엔진 $N$대};
  \draw[rounded corners=2pt,draw=cblue,fill=white] (6.2,-2.7) rectangle (14.2,1.2);
  \node[anchor=north west,font=\small\bfseries] at (6.3,1.15) {\code{bombe\_engine}};
  \node[blko,minimum width=21mm] (cfg) at (7.5,-0.1) {설정 레지스터\\{\scriptsize 메뉴, 로터 순서}};
  \node[blkb,minimum width=21mm] (tbl) at (10.2,-0.1) {치환표 계산\\{\scriptsize 파이프라인}};
  \node[reg,minimum width=20mm] (ram) at (12.8,-0.1) {표 메모리\\{\scriptsize 분산 RAM}};
  \node[blkb,minimum width=21mm] (spr) at (10.2,-1.8) {전파\\{\scriptsize 676비트 상태}};
  \node[blkr,minimum width=21mm] (fifo) at (7.5,-1.8) {정지 FIFO};
  \draw[->] (cfg) -- (tbl); \draw[->] (tbl) -- (ram);
  \draw[->] (ram.south) |- (spr.east);
  \draw[->] (spr) -- (fifo);
  \draw[<->] (host) -- (br);
  \draw[<->] (br) -- node[above,font=\scriptsize,text=cgray] {레지스터 버스} (6.2,-0.1);
\end{tikzpicture}
\caption{FPGA 엔진의 블록 수준 구성.}
\label{fig:g-engine-blocks}
\end{figure}

엔진이 옳다는 근거는 3부와 같은 방식이다. 시뮬레이션에서 엔진의 정지 목록을 Rust 코어의 정지 목록과 순서까지
비교한다.

\begin{todobox}{\code{bombe\_engine} 해설}
\begin{itemize}
\item 구동 시작 때 고른 로터의 배선표를 레지스터로 옮기는 부분과 그 이유
\item 안쪽 치환 파이프라인(14단)의 코드와 단별 타이밍 그림
\item 오른쪽 로터 26위치 치환표의 병렬 레인(5단)과 분산 RAM에 쓰는 과정
\item 메뉴의 선마다 표의 사본을 두는 이유(분산 RAM의 읽기 단자)
\item 케이블 단위 전파 회로의 코드를 한 줄씩, 구조 그림과 함께
\item 고정점 판정을 레지스터에 담긴 값으로 하는 이유와 조기 종료
\item 엔진의 상태 기계 전체와 위치당 클럭 수
\item 정지 FIFO와 가득 찼을 때의 동작
\end{itemize}
\end{todobox}

\begin{todobox}{\code{bombe\_array}, \code{uart\_bridge}, \code{axil\_bridge} 해설}
\begin{itemize}
\item 레지스터 버스: 주소로 엔진을 고르는 방식, 요청과 응답의 타이밍
\item UART 브리지: 직렬 통신의 원리(일반적인 예로 UART 송수신기)와 읽기·쓰기 명령의 상태 기계
\item AXI4-Lite 브리지: 보드의 처리 장치가 FPGA 안의 레지스터를 읽고 쓰는 방식
\item 보드의 최상위 모듈과 제약 파일이 하는 일
\end{itemize}
\end{todobox}

\begin{todobox}{배치·배선 이후}
확정된 자원 사용량과 타이밍, 엔진 수에 따른 변화, 보드에서의 실측 속도. 전시물로서 내부 상태(676가닥의 선,
드럼 위치, 정지)를 눈에 보이게 하는 표시 장치의 구조.
\end{todobox}
